% Options for packages loaded elsewhere
\PassOptionsToPackage{unicode}{hyperref}
\PassOptionsToPackage{hyphens}{url}
\documentclass[
  11pt,
]{article}
\usepackage{xcolor}
\usepackage[margin=2.5cm]{geometry}
\usepackage{amsmath,amssymb}
\setcounter{secnumdepth}{-\maxdimen} % remove section numbering
\usepackage{iftex}
\ifPDFTeX
  \usepackage[T1]{fontenc}
  \usepackage[utf8]{inputenc}
  \usepackage{textcomp} % provide euro and other symbols
\else % if luatex or xetex
  \usepackage{unicode-math} % this also loads fontspec
  \defaultfontfeatures{Scale=MatchLowercase}
  \defaultfontfeatures[\rmfamily]{Ligatures=TeX,Scale=1}
\fi
\usepackage{lmodern}
\ifPDFTeX\else
  % xetex/luatex font selection
\fi
% Use upquote if available, for straight quotes in verbatim environments
\IfFileExists{upquote.sty}{\usepackage{upquote}}{}
\IfFileExists{microtype.sty}{% use microtype if available
  \usepackage[]{microtype}
  \UseMicrotypeSet[protrusion]{basicmath} % disable protrusion for tt fonts
}{}
\makeatletter
\@ifundefined{KOMAClassName}{% if non-KOMA class
  \IfFileExists{parskip.sty}{%
    \usepackage{parskip}
  }{% else
    \setlength{\parindent}{0pt}
    \setlength{\parskip}{6pt plus 2pt minus 1pt}}
}{% if KOMA class
  \KOMAoptions{parskip=half}}
\makeatother
\usepackage{longtable,booktabs,array}
\usepackage{caption}
\captionsetup[table]{skip=6pt}
\newcounter{none} % for unnumbered tables
\usepackage{calc} % for calculating minipage widths
% Correct order of tables after \paragraph or \subparagraph
\usepackage{etoolbox}
\makeatletter
\patchcmd\longtable{\par}{\if@noskipsec\mbox{}\fi\par}{}{}
\makeatother
\setlength{\emergencystretch}{3em} % prevent overfull lines
\providecommand{\tightlist}{%
  \setlength{\itemsep}{0pt}\setlength{\parskip}{0pt}}
\usepackage[]{natbib}
\bibliographystyle{plainnat}
\usepackage{seqsplit}
\usepackage{bookmark}
\IfFileExists{xurl.sty}{\usepackage{xurl}}{} % add URL line breaks if available
\urlstyle{same}
% fallback for those not using the hyperref driver hyperxmp:
\makeatletter
\@ifundefined{xmpquote}{\newcommand{\xmpquote}[1]{#1}}{}
\makeatother
\hypersetup{
  hidelinks,
  pdfcreator={LaTeX via pandoc}}

\author{}
\date{}

\begin{document}

\section{Largest Lyapunov exponent estimators on the standard map: a
preregistered numerical
study}\label{largest-lyapunov-exponent-estimators-on-the-standard-map-a-preregistered-numerical-study}

\subsection{Abstract}\label{abstract}

We report a preregistered numerical study of largest-Lyapunov-exponent
estimators on the standard map. Scope: one panel from a single draw (a
single seed), run in Float64, with five nonzero kick values and one
start per kick value in the reproduction test (H4). Section 10 gives the
limitations. The panel had 883 orbits: 13 exact-exponent controls,
orbits at K = 0, and orbits at five nonzero kick values. The decision
rules were frozen before any data existed. The claim of asymptotic
equivalence (H1) was falsified. The failure is carried by sub-rule H1c
alone, with 162 of 209 eligible orbits passing, a pass fraction of
0.775119617225 against a limit of 0.90. Sub-rules H1a and H1b were met
in full. H1 is evaluated in Float64 at the horizon N\_1 = 1e7, and H1c
compares L2 run in wrapped coordinates with L5 in lifted coordinates
(limitation 4, Section 10). The claim of finite-time inequivalence (H2)
was not supported (NEGATIVE). Over the 355 orbits of the preregistered
set W the median of R(n\_H2) was 2.23588035691e-5 against a threshold of
0.10. Set W is defined in Section 2.2. The regime dependence of the
convergence time (H3) was supported, with a ratio of medians of
11213.7767857. The reproduction of the reference values of Sala and
Artuso \citep{SalaArtuso2015} (H4) did not reproduce at K = pi/2, where
we obtained 0.256066843848 against 0.298. It passed at K = 2 pi, where
we obtained 1.17018723872 against 1.172. The label rule H5 was supported
under its three sub-rules (H5a, H5b, H5c). The ladder test forced 35 of
the 92 ladder-subset labels to unresolved; this is reported as
exploratory item D3. Cost and sensitivity are reported without a
decision rule. The run used package version 0.0.1, before the tagged
release v0.1.0. The code is released as Chaoslab v0.1.1 (TOOL), a tagged
and documented release with the same computing code (Section 11).

\subsection{1. Introduction}\label{introduction}

This paper reports a preregistered numerical study of
largest-Lyapunov-exponent estimators on the standard map. The study
tests two claims of the project, C1 (asymptotic equivalence) and C2
(finite-time inequivalence). The claim texts below are quoted from the
project file \texttt{\_config/research-question.md}. They are claims of
the study. They are not statements supported by the literature cited in
this paper.

Claim C1, as stated in \texttt{\_config/research-question.md}:
``Benettin, QR (Gram-Schmidt / Householder), stabilized SVD and
well-conditioned tangent-map methods agree on lambda\_1 when run long
enough at adequate precision.'' Claim C2, as stated in
\texttt{\_config/research-question.md}: ``Finite-time LE estimators
differ substantially over realistic budgets, most of all in weak chaos,
sticky layers and near separatrices.'' Here LE stands for Lyapunov
exponent and lambda\_1 for the largest exponent. The tests of C2 in this
study run over the preregistered set W of Section 2.2, whose composition
is shown in Tables 6 and 8.

The study also attempts to reproduce two published values. Sala and
Artuso give 0.298 at K = pi/2 and 1.172 at K = 2 pi
\citep{SalaArtuso2015}.

Two published comparators serve for the chaotic-sea exponent. These are
the fit of Shevchenko \citep{Shevchenko2014} and the line ln(K/2) of
Chirikov \citep{Chirikov1979}. The study uses both as comparators that
never pass or fail.

Lyapunov spectra are usually computed by renormalising tangent vectors
and orthogonalising them by Gram-Schmidt
\citep{Benettin1980a, Benettin1980b, ShimadaNagashima1979}.
Gram-Schmidt, Householder and continuous QR variants were compared on a
Toda chain; the discrete variants differed by much less than the
fluctuations of the finite-time values, and continuous QR took about 40
times the CPU time of Gram-Schmidt \citep{Geist1990}. Singular-value
methods take the singular values of a product of Jacobians without
forming the product \citep{Stewart1997}, and their theory holds for
systems with stable and distinct exponents \citep{DieciElia2006}.

Finite-time effects are known in particular settings. For standard-like
maps the finite-time exponent carries an initial-vector transient that
fades like 1/n \citep{SalaArtuso2015}. For the two-trajectory method on
flows, the renormalisation time and the initial distance can produce a
spurious largest exponent \citep{DubeibeBermudezAlmanza2014}. Sticky
orbits need long runs: a finite-time exponent in an extreme sticky zone
of the standard map had not converged after 5 x 10\^{}7 iterations,
against about 10\^{}7 in the chaotic sea
\citep{HarsoulaKaramanosContopoulos2019}, and one orbit at k = 0.975
looks confined for 10\^{}6 steps and escapes after 2.3 x 10\^{}9
\citep{Meiss2015}. Round-off induced divergence of orbits has been
related to the reversibility error in discrete-time systems
\citep{FarandaMestreTurchetti2012}. This study compares the estimators
above on one map, over five orbit regimes, with a precision ladder on a
subset of the orbits.

This paper reports the study as preregistered. It adds no new claim,
softens no result and offers no reinterpretation. Section 2 gives the
design, the notation and the run identity. Sections 3 to 7 give the
confirmatory results for hypotheses H1 to H5. Exploratory subsections
inside those sections are marked EXPLORATORY: D1 and D2 in Section 3.5,
D4 in Section 6.4, and D3 and D5 in Section 7.3. Section 8 gives cost
and sensitivity. Section 9 gives the exploratory items D6 and D7.
Sections 10 to 12 give limitations, availability and the conclusion.

\subsection{2. Methods and
preregistration}\label{methods-and-preregistration}

\subsubsection{2.1 Preregistration}\label{preregistration}

The decision rules were frozen on 2026-10-17, before any data existed.
Amendments A1 to A6 were all frozen before the first study data file.
The rules are applied as frozen, with the deviations listed in Section
2.5. The results document also lists 14 readings where the frozen text
is silent (Section 2.5). Amendments A1 to A4 are described here by
pointer to the frozen text. A1 sets the hash rule, the T predicate, the
screening direction, the windowed estimate, the Newton periodic-point
search, the subset caps and the machine-readable panel inputs {[}prereg:
Amendment A1, items 1 to 7{]}. A2 handles a ghost pair of periodic
points at K = 2 pi and is superseded by A3 {[}prereg: Amendment A2,
Amendment A3{]}. A3 sets the acceptance test of the periodic-point
search {[}prereg: Amendment A3{]}. A4 adds the machine-readable run
block, the setup check, the cost cap and pilot, the backend agreement
and the scope of the control bound {[}prereg: Amendment A4, items 1 to
5{]}. The content of A5 and A6 is given in Section 2.5. The
preregistered value types were FALSIFIES (C1, C2), REPRODUCES
(\citep{SalaArtuso2015}), TOOL and NEGATIVE allowed. Claim C3 is out of
scope. TOOL claim: the library Chaoslab, release v0.1.1. Every method
the run used has catalog status \texttt{verified} (preregistration,
Methods). Release v0.1.1 is tagged and carries user documentation, and
its package code differs from the analysis commit only in docstrings and
in a file-descriptor conversion that is a no-op on Julia 1.13.2, the
Julia version of the run (Section 11). The preregistration header is
identified by the SHA-256 value in Table 1. The location of the frozen
text is given in Section 11.

\subsubsection{2.2 Notation and
definitions}\label{notation-and-definitions}

This subsection restates the notation of the study design. Each item has
a pointer to the frozen text. The definition of the standard map and of
its tangent map is in \texttt{\_config/research-question.md} and is not
repeated here (pointer only).

\begin{itemize}
\tightlist
\item
  Variables and horizons. The map variables are (x, p). The H4 start is
  x0 = 1e-3 and p0 = -1e-3. K is the kick parameter. The horizons are
  N\_1 = 1e7 (all orbits), N\_long = 1e8 (designated subset) and N\_B =
  1e6 (precision ladder). N is the horizon of the series used; R(n) and
  n\_conv use the horizon-N\_1 series of every orbit.
\item
  Finite-time exponent. lambda\_m(n) is the finite-time largest Lyapunov
  exponent of method m after n map steps, in natural log per step, with
  the Euclidean 2-norm on the lifted (x, p) tangent vector. L9 uses the
  largest finite-time singular value of the product. lambda\_ref of an
  orbit is lambda\_L5 at the longest Float64 horizon of that orbit.
\item
  Methods. L1 is Benettin. L2 is the two-trajectory method, an
  approximation of the largest exponent. L3 is CGS, L4 is MGS, L4b is
  RGS, L5 is Householder (the default reference), L5b is the H2 diagonal
  iteration path, L6 is Givens, and L9 is the stabilized SVD (d = 2).
  Their algorithm details are in the study design and the repository,
  not here.
\item
  Method sets. Q = \{L3, L4, L4b, L5, L5b, L6\}. The variational set is
  V = \{L1\} + Q + \{L9\}. The expression ``\{L1\} + Q'' means L1
  together with the members of Q.
\item
  Parameters. tau is the fixed reorthonormalization interval of the Q
  methods (primary tau = 1). eps is the perturbation size of L2 (primary
  eps\_0 = 1e-8). The primary initial deviation direction is the angle
  pi/4 from the x axis. On the designated subset four more ``direction
  offsets'' are used: 0.3, 0.3 + pi/5, 0.3 + 2 pi/5 and 0.3 + 3 pi/5.
  The precisions are Float64, Double64 and BigFloat.
\item
  Spreads and ratios. Delta\_QB(n) is the maximum minus the minimum of
  lambda\_m(n) over \{L1\} + Q. Delta\_V(n) is the same over V. R(n) =
  Delta\_V(n) / lambda\_ref (H2). n\_conv is defined in H3 (Section
  5.1); the series are stored at recorded values of n, so n\_conv takes
  recorded values of n.~
\item
  Design classes. T (KAM torus), R (resonant), C (chaotic sea), S
  (sticky) and X (near a separatrix or a hyperbolic fixed point). Design
  class is the a-priori intended class of an orbit. A cell is a K value
  times a design class.
\item
  Controls. The controls are the exact periodic points of the map and
  one hyperbolic linear map (Y5), whose exponents are known
  analytically. ``Info'' controls are defined in Section 2.5.
\item
  Sets. E\_ch is the set of design class C orbits at K \textgreater{} 0
  plus the controls with positive analytic exponent (H1). The weak-chaos
  set W is defined below. The strong-chaos set S is the class C orbits
  at K = 2 pi and K = 10 (H3). S\_D is the designated subset and S\_L is
  the ladder subset, nested in the panel; S\_L carries the Double64 and
  BigFloat runs. The symbol S names three things in this paper: the
  panel class S (sticky), the strong-chaos set S of H3 and, as S\_L, the
  ladder subset. The context says which.
\item
  Definition of W (preregistered). W = orbits with instability label
  ``positive growth supported'' in the design classes sticky or
  near-separatrix at any K \textgreater{} 0, plus chaotic-sea orbits at
  K \textless= K\_c.~ The name ``weak-chaos set'' is the
  preregistration's. By this definition W contains class S and class X
  orbits at the large K values as well, and Tables 6 and 8 show its
  cells. The name does not mean that every member of W is in weak chaos.
\item
  Labels. Each orbit carries two labels. The instability label takes the
  values ``positive growth supported'', ``none detected on the tested
  horizon'' and ``unresolved''. The structure label takes the values
  ``periodic'', ``quasiperiodic'', ``irregular'' and ``unresolved''.
  Four disagreement tests force a label to ``unresolved'': the precision
  ladder test (on S\_L only), the estimators test (any failure of H1a or
  H1b on that orbit), the direction test (on the designated subset only)
  and the diagnostics test (a pairing, determinant or orthogonality
  residual above its tolerance for L5). A label is ``forced'' when one
  of these tests sets it to ``unresolved''. The thresholds of the label
  rules are in the frozen text.
\item
  Drift and kept orbits. An orbit is flagged ``drift'' when its lifted
  \textbar p\textbar{} exceeds 1e6 at any recorded time. Drift orbits
  are removed from E\_ch, W and S. ``Kept'' orbits are the orbits not
  dropped as drift or excluded.
\item
  Bootstrap. The bootstrap intervals are percentile intervals at the 95
  percent level with 10\^{}4 resamples. Orbits are resampled with
  replacement, and W and S are resampled independently.
\end{itemize}

\subsubsection{2.3 Panel and sets}\label{panel-and-sets}

The study compares estimators on a panel of 883 orbits. The panel table
{[}T1{]} gives its cells: 13 controls in the cell of K index 0; 40
orbits each in classes R and T at K index 1; and the cells at K index 2
to 6. At K index 2 the classes are C, R, S, T and X with 40, 40, 25, 40
and 40 orbits. At K index 3 and 4 the classes are C, R, S, T and X with
40 orbits each. At K index 5 the classes are C, S, T and X with 40, 5,
40 and 40 orbits. At K index 6 the classes are C and X with 40 orbits
each. Cells that are short, empty or undefined are listed in D6 (Section
9). K index 0 holds the exact-exponent controls. K index 1 is K = 0. K
index 2 to 6 are K = 0.5, 0.971635406 (K\_c, the breakup value of the
golden-mean KAM torus, as given in
\citep{FrahmShepelyansky2010, CarlettiLaskar2000} and fixed in the
preregistration), 1.57079632679 (pi/2), 6.28318530718 (2 pi) and 10.
``Class S'' is a panel class. It is not the ``strong-chaos set S'' of
H3, which is the class C orbits at K = 2 pi and K = 10. One orbit, orbit
881, is dropped as drift. Under the preregistration, drift orbits are
removed from E\_ch, W and S. H5b is defined on all panel orbits, and the
results document gives 883 for it. The sets used by the decision rules
have these sizes: E\_ch 209, W 355, S 80 and S\_L 92. H5b uses 883
orbits and H5c uses 334. The program reads S as the class C orbits at K
= 2 pi and K = 10 without a label condition. It reads E\_ch as class C
at K \textgreater{} 0 (plus the controls, Section 3.1). For W this
paper, the results document and the findings file use the preregistered
definition of Section 2.2. The analysis program prints a shorter phrase
for W in its list of readings, class C at 0 \textless{} K \textless=
K\_c {[}findings E, reading 8{]}. That phrase is not the frozen
definition and does not describe the cells of Tables 6 and 8, which also
contain classes S and X (limitation 6, Section 10). The estimator
methods L1 to L9 and Q are defined in the study design (Section 2.2).
The repository (Section 11) holds the details.

\subsubsection{2.4 Run identity}\label{run-identity}

All reported numbers come from one fully automatic analysis run over the
stored series. The run completed with 1363 chunks, all with status ok.
Table 1 identifies the run.

Table 1. Run identity.

{\def\LTcaptype{none} % do not increment counter
\begin{longtable}[]{@{}
  >{\raggedright\arraybackslash}p{(\linewidth - 4\tabcolsep) * \real{0.3333}}
  >{\raggedright\arraybackslash}p{(\linewidth - 4\tabcolsep) * \real{0.3333}}
  >{\raggedright\arraybackslash}p{(\linewidth - 4\tabcolsep) * \real{0.3333}}@{}}
\toprule\noalign{}
\begin{minipage}[b]{\linewidth}\raggedright
Item
\end{minipage} & \begin{minipage}[b]{\linewidth}\raggedright
Value
\end{minipage} & \begin{minipage}[b]{\linewidth}\raggedright
Pointer
\end{minipage} \\
\midrule\noalign{}
\endhead
\bottomrule\noalign{}
\endlastfoot
Git commit (tables header and findings) &
\seqsplit{d31b655d0c33e2f2f1a67bbbdeba9c8c8a3fd036} & {[}findings A, table{]}, {[}T
Header: git commit{]} \\
Git sha recorded by the run manifest &
\seqsplit{3e405f3324f8d25162580d3380db26c877d7ca63} & {[}manifest: git.sha{]} \\
Module version & 0.0.1 & {[}T Header: module version{]}, {[}manifest:
versions.module{]} \\
Panel SHA-256 &
\seqsplit{881874d26aa8883caa0545331ba167a3f2c9e3c249cdb6ba32ad39c98ad348d5} &
{[}findings A, table{]}, {[}manifest: panel.sha256\_end{]} \\
Preregistration header SHA-256 &
\seqsplit{62011ba5af514da4f1238aef9ae6cb35ca21093c4343b71dd696df360b9994bf} &
{[}findings A, table{]}, {[}manifest: prereg.header{]} \\
Manifest SHA-256 &
\seqsplit{ae6ec80eb08fbd6e98aeef1fbf840a567a03102ab943f2ae0aaf2b3e367f191d} &
{[}findings A, table{]} \\
Parameter file SHA-256 &
\seqsplit{a544130bc56dc66d8b1f7ec5d1e54197306f14dc19c7cd22fdbd94e4534b7282} &
{[}findings A, table{]} \\
Bootstrap seed & 20261012 & {[}findings A, table{]}, {[}params:
bootstrap seed in Header{]} \\
Run status & complete, 1363 chunks, all with status ok & {[}manifest:
status{]}, {[}manifest: plan\_counts.chunks{]}, {[}manifest:
chunks{[}{]}.status{]} \\
\end{longtable}
}

The tables header and the findings file give the first git commit. The
run manifest records the second. The run commit is an ancestor of the
analysis commit, and between them only the analysis code, its figure
extension, the command-line entry point and the project file change. The
stored series therefore come from the run commit and the tables and
figures from the analysis commit, with the same series code at both
(limitation 8, Section 10).

\subsubsection{2.5 Deviations from the frozen text and
readings}\label{deviations-from-the-frozen-text-and-readings}

The frozen text was applied with the deviations that the program states.

\begin{enumerate}
\def\labelenumi{\arabic{enumi}.}
\tightlist
\item
  L2 runs in wrapped coordinates (Amendment A5 item 1). H1c compares the
  wrapped L2 value with the lifted L5 value of the same orbit.
\item
  The tau sweep records an orthonormalisation breakdown as a result, not
  as a failure (A5 item 2).
\item
  The control bound applies to controls 871, 873, 875, 877, 879 and 883
  (A5 item 3). Controls 878, 880, 881 and 882 are reported as info and
  are not held to the bound.
\item
  The control bound of L9 is 5.0e-8 and that of the other methods is
  1.0e-8 (A6).
\item
  The manifest records no overrides.
\end{enumerate}

The results document lists 14 readings where the frozen text is silent
{[}findings E, Readings 1 to 14; T Readings{]}. It names two that bear
directly on the verdicts. Reading 5: H5a is read from the per-precision
labels before the ladder test acts, and H5 is supported only if H5a, H5b
and H5c all pass. Reading 8: the definitions of S, W and E\_ch as the
program reads them (Section 2.3). The other twelve readings are listed
in the tables file, section Readings, and are not described in this
paper.

\subsection{3. Results: H1 (claim C1),
FALSIFIED}\label{results-h1-claim-c1-falsified}

\subsubsection{3.1 Preregistered rule}\label{preregistered-rule}

The rule applies on E\_ch, which is class C orbits at K \textgreater{} 0
plus the exact-exponent controls. The horizon is N\_1 = 1e7. An orbit
passes H1a if the spread over \{L1\} + Q is at most 1e-8. It passes H1b
if \textbar lambda\_L9 - lambda\_L5\textbar{} is at most 1e-6. It passes
H1c if \textbar lambda\_L2 - lambda\_L5\textbar{} is at most 5e-3 at
eps\_0 = 1e-8. H1 is SUPPORTED only if the pass fraction is at least
0.98 for H1a and H1b and at least 0.90 for H1c. It is FALSIFIED if any
sub-rule is below its limit.

\subsubsection{3.2 Measured}\label{measured}

There were 209 eligible orbits. The verdict as printed is falsified.
Table 2 gives the three sub-rules.

Table 2. H1 sub-rules on E\_ch.

{\def\LTcaptype{none} % do not increment counter
\begin{longtable}[]{@{}
  >{\raggedright\arraybackslash}p{(\linewidth - 10\tabcolsep) * \real{0.1667}}
  >{\raggedright\arraybackslash}p{(\linewidth - 10\tabcolsep) * \real{0.1667}}
  >{\raggedright\arraybackslash}p{(\linewidth - 10\tabcolsep) * \real{0.1667}}
  >{\raggedright\arraybackslash}p{(\linewidth - 10\tabcolsep) * \real{0.1667}}
  >{\raggedright\arraybackslash}p{(\linewidth - 10\tabcolsep) * \real{0.1667}}
  >{\raggedright\arraybackslash}p{(\linewidth - 10\tabcolsep) * \real{0.1667}}@{}}
\toprule\noalign{}
\begin{minipage}[b]{\linewidth}\raggedright
Sub-rule
\end{minipage} & \begin{minipage}[b]{\linewidth}\raggedright
Limit
\end{minipage} & \begin{minipage}[b]{\linewidth}\raggedright
Passing orbits
\end{minipage} & \begin{minipage}[b]{\linewidth}\raggedright
Fraction
\end{minipage} & \begin{minipage}[b]{\linewidth}\raggedright
Result
\end{minipage} & \begin{minipage}[b]{\linewidth}\raggedright
Pointer
\end{minipage} \\
\midrule\noalign{}
\endhead
\bottomrule\noalign{}
\endlastfoot
H1a & 0.98 & 209 of 209 & 1.0 & pass & {[}T2, header line; params:
h1a\_fraction{]} \\
H1b & 0.98 & 209 of 209 & 1.0 & pass & {[}T2, header line; params:
h1b\_fraction{]} \\
H1c & 0.90 & 162 of 209 & 0.775119617225 & fail & {[}T2, header line;
params: h1c\_fraction{]} \\
\end{longtable}
}

H1 is falsified. The H1c pass fraction was 162 of 209, which is
0.775119617225, against the limit of 0.90. The shortfall is carried by
sub-rule H1c alone. H1a and H1b were met in full.

Table 3 gives H1c by cell.

Table 3. H1c passing orbits by cell.

{\def\LTcaptype{none} % do not increment counter
\begin{longtable}[]{@{}llll@{}}
\toprule\noalign{}
Cell & Orbits & H1c pass & Pointer \\
\midrule\noalign{}
\endhead
\bottomrule\noalign{}
\endlastfoot
controls & 9 & 9 & {[}T2, row controls{]} \\
K index 2 & 40 & 27 & {[}T2, row K index 2{]} \\
K index 3 & 40 & 13 & {[}T2, row K index 3{]} \\
K index 4 & 40 & 33 & {[}T2, row K index 4{]} \\
K index 5 & 40 & 40 & {[}T2, row K index 5{]} \\
K index 6 & 40 & 40 & {[}T2, row K index 6{]} \\
\end{longtable}
}

The 47 failing orbits sit in three K cells. At K index 2, 27 of 40
orbits pass. At K index 3, 13 of 40 pass. At K index 4, 33 of 40 pass. K
index 5 and 6 pass 40 of 40, and the 9 controls pass. H1a and H1b pass
in every cell. The ids of the failing orbits and the extreme deviations
are in {[}T2, Failing orbits{]}. Every failing orbit has
\textbar lambda\_L2 - lambda\_L5\textbar{} above 0.005. The smallest is
0.00548462883667 and the largest is 0.0353748064771. On the failing
orbits the largest Delta\_QB is 9.66310365058e-14 and the largest
\textbar lambda\_L9 - lambda\_L5\textbar{} is 3.58331395889e-7. None of
them breaks H1a or H1b.

\subsubsection{3.3 Exact-exponent
controls}\label{exact-exponent-controls}

The controls under the bound are 871, 873, 875, 877, 879 and 883. Table
4 gives the largest deviation over the methods.

Table 4. Exact-exponent controls under the bound.

{\def\LTcaptype{none} % do not increment counter
\begin{longtable}[]{@{}
  >{\raggedright\arraybackslash}p{(\linewidth - 8\tabcolsep) * \real{0.2000}}
  >{\raggedright\arraybackslash}p{(\linewidth - 8\tabcolsep) * \real{0.2000}}
  >{\raggedright\arraybackslash}p{(\linewidth - 8\tabcolsep) * \real{0.2000}}
  >{\raggedright\arraybackslash}p{(\linewidth - 8\tabcolsep) * \real{0.2000}}
  >{\raggedright\arraybackslash}p{(\linewidth - 8\tabcolsep) * \real{0.2000}}@{}}
\toprule\noalign{}
\begin{minipage}[b]{\linewidth}\raggedright
Control
\end{minipage} & \begin{minipage}[b]{\linewidth}\raggedright
Scope
\end{minipage} & \begin{minipage}[b]{\linewidth}\raggedright
Largest deviation
\end{minipage} & \begin{minipage}[b]{\linewidth}\raggedright
L9 deviation
\end{minipage} & \begin{minipage}[b]{\linewidth}\raggedright
Pointer
\end{minipage} \\
\midrule\noalign{}
\endhead
\bottomrule\noalign{}
\endlastfoot
871 & bound & 5.26802590439e-9 & 5.26802579337e-9 & {[}T2, controls,
871{]} \\
873 & bound & 2.48736364838e-9 & 8.326783706990001e-12 & {[}T2,
controls, 873{]} \\
875 & bound & 5.16177345311e-9 & 1.82781789704e-9 & {[}T2, controls,
875{]} \\
877 & bound & 1.79536172595e-8 & 1.79536172595e-8 & {[}T2, controls,
877{]} \\
879 & bound & 2.28260139679e-8 & 2.28260139679e-8 & {[}T2, controls,
879{]} \\
883 & bound & 1.11022302463e-16 & 1.11022302463e-16 & {[}T2, controls,
883{]} \\
\end{longtable}
}

The largest deviations lie within the bounds of Amendment A6. For
controls 877 and 879 the largest deviation is the L9 deviation. It is
above 1e-8 and below the L9 control bound of 5.0e-8. The table does not
give the other methods' deviations on these two controls separately.

Accounting for the control ids. The panel table counts 13 controls in
the cell of K index 0, and the H1 cell table counts 9 controls in
E\_ch.~The controls table lists 10 ids: the six of Table 4 and the four
info controls 878, 880, 881 and 882 (D2). The other 3 controls, orbits
872, 874 and 876, start at the elliptic fixed point (pi, 0) at K = 0.5,
K\_c and pi/2 and have analytic exponent 0. E\_ch and the controls table
take only controls with positive analytic exponent, so these three
appear in neither; they serve the setup check of the Double64 and
BigFloat runs and enter H5b with every panel orbit. Orbit 881 is an info
control and is the orbit dropped as drift (Section 2.3).

\subsubsection{3.4 Reading}\label{reading}

The preregistered rule for C1 is not met. H1 is reported as falsified.
The falsification rests on sub-rule H1c alone. It was evaluated in
Float64 at N\_1 = 1e7 on one panel, and H1c compares L2 in wrapped
coordinates with L5 in lifted coordinates (Section 2.5, item 1;
limitation 4 in Section 10). The tables do not show whether the
coordinate choice accounts for any part of the H1c result.

\subsubsection{3.5 Exploratory items D1 and D2
(EXPLORATORY)}\label{exploratory-items-d1-and-d2-exploratory}

Everything in this subsection is EXPLORATORY. No frozen rule applies and
no cause is claimed.

D1 (EXPLORATORY): where the H1c failures sit. The 47 H1c failures are in
three K cells. K index 2 has 13 of 40 failing (27 pass). K index 3 has
27 of 40 (13 pass). K index 4 has 7 of 40 (33 pass). K index 5 and 6
have no failures (40 of 40 pass each) and the 9 controls all pass. The
failing ids fall in three ranges: 13 ids between 166 and 196, 27 between
346 and 385, and 7 between 549 and 580. These counts match the per-cell
failure counts, but the tables do not state the orbit-to-cell mapping.
The tables do not show why the failures sit at the three smaller K.

D2 (EXPLORATORY): the info controls. Controls 878, 880, 881 and 882 are
reported as info. Their largest deviations are 0.222903869545,
0.443364672685, 0.650205817015 and 0.103506221066. These are far above
the bound that applies to the other six controls. Amendment A5 item 3
designates these four as the controls reported as info; the results
document describes them as the controls whose binary64 start is not an
exact periodic point. That description is not tested here. The tables do
not say more.

\subsection{4. Results: H2 (claim C2), NOT SUPPORTED
(NEGATIVE)}\label{results-h2-claim-c2-not-supported-negative}

\subsubsection{4.1 Preregistered rule}\label{preregistered-rule-1}

R(n) = Delta\_V(n) / lambda\_ref is evaluated at n\_H2 = 1e4 steps over
the set W defined in Section 2.2. The preregistration names W the
weak-chaos set. W = orbits with instability label ``positive growth
supported'' in the design classes sticky or near-separatrix at any K
\textgreater{} 0, plus chaotic-sea orbits at K \textless= K\_c.~ H2 is
SUPPORTED if the median of R(n\_H2) is at least 0.10 and the lower end
of its 95 percent bootstrap interval (10\^{}4 resamples of orbits) is at
least 0.05. Otherwise the null is reported.

\subsubsection{4.2 Measured}\label{measured-1}

W has 355 orbits. The median of R(n\_H2) is 2.23588035691e-5. The 95
percent bootstrap interval is {[}1.9257235885e-5, 2.31226443693e-5{]}.
The verdict as printed is not supported. Both conditions fail. The
median is below 0.10 and the lower end is below 0.05.

Table 5. H2 against the preregistered thresholds.

{\def\LTcaptype{none} % do not increment counter
\begin{longtable}[]{@{}
  >{\raggedright\arraybackslash}p{(\linewidth - 8\tabcolsep) * \real{0.2000}}
  >{\raggedright\arraybackslash}p{(\linewidth - 8\tabcolsep) * \real{0.2000}}
  >{\raggedright\arraybackslash}p{(\linewidth - 8\tabcolsep) * \real{0.2000}}
  >{\raggedright\arraybackslash}p{(\linewidth - 8\tabcolsep) * \real{0.2000}}
  >{\raggedright\arraybackslash}p{(\linewidth - 8\tabcolsep) * \real{0.2000}}@{}}
\toprule\noalign{}
\begin{minipage}[b]{\linewidth}\raggedright
Quantity
\end{minipage} & \begin{minipage}[b]{\linewidth}\raggedright
Measured
\end{minipage} & \begin{minipage}[b]{\linewidth}\raggedright
Threshold
\end{minipage} & \begin{minipage}[b]{\linewidth}\raggedright
Result
\end{minipage} & \begin{minipage}[b]{\linewidth}\raggedright
Pointer
\end{minipage} \\
\midrule\noalign{}
\endhead
\bottomrule\noalign{}
\endlastfoot
Median of R(n\_H2) over W (355 orbits) & 2.23588035691e-5 & at least
0.10 & fail & {[}T3, header line; params: h2\_median{]} \\
Lower end of the 95 percent interval & 1.9257235885e-5 & at least 0.05 &
fail & {[}T3, header line; params: h2\_interval{]} \\
\end{longtable}
}

This is a null (NEGATIVE) result. It has the same standing as a pass. In
W, the relative spread among the methods of V at 1e4 steps is far below
the preregistered level. The frozen rule decides on the pooled W. Table
6 gives the cells of W; the cell results are descriptive. W is composed
of class C at K index 2 and 3, class S at K index 2 to 5 and class X at
K index 2 to 6.

Table 6. H2: R(n\_H2) by cell of W (descriptive).

{\def\LTcaptype{none} % do not increment counter
\begin{longtable}[]{@{}
  >{\raggedright\arraybackslash}p{(\linewidth - 14\tabcolsep) * \real{0.1250}}
  >{\raggedright\arraybackslash}p{(\linewidth - 14\tabcolsep) * \real{0.1250}}
  >{\raggedright\arraybackslash}p{(\linewidth - 14\tabcolsep) * \real{0.1250}}
  >{\raggedright\arraybackslash}p{(\linewidth - 14\tabcolsep) * \real{0.1250}}
  >{\raggedright\arraybackslash}p{(\linewidth - 14\tabcolsep) * \real{0.1250}}
  >{\raggedright\arraybackslash}p{(\linewidth - 14\tabcolsep) * \real{0.1250}}
  >{\raggedright\arraybackslash}p{(\linewidth - 14\tabcolsep) * \real{0.1250}}
  >{\raggedright\arraybackslash}p{(\linewidth - 14\tabcolsep) * \real{0.1250}}@{}}
\toprule\noalign{}
\begin{minipage}[b]{\linewidth}\raggedright
Class
\end{minipage} & \begin{minipage}[b]{\linewidth}\raggedright
K index
\end{minipage} & \begin{minipage}[b]{\linewidth}\raggedright
W orbits
\end{minipage} & \begin{minipage}[b]{\linewidth}\raggedright
Median
\end{minipage} & \begin{minipage}[b]{\linewidth}\raggedright
Lo
\end{minipage} & \begin{minipage}[b]{\linewidth}\raggedright
Hi
\end{minipage} & \begin{minipage}[b]{\linewidth}\raggedright
Note
\end{minipage} & \begin{minipage}[b]{\linewidth}\raggedright
Pointer
\end{minipage} \\
\midrule\noalign{}
\endhead
\bottomrule\noalign{}
\endlastfoot
C & 2 & 36 & 0.000685409790935 & 1.21772360607e-5 & 0.00158551651977 & &
{[}T3, row C 2{]} \\
S & 2 & 21 & 5.42291035617e-5 & 3.17071908742e-6 & 0.00145378914565 & &
{[}T3, row S 2{]} \\
X & 2 & 36 & 5.82402750446e-12 & 7.73045459959e-13 & 3.49257079589e-11 &
& {[}T3, row X 2{]} \\
C & 3 & 34 & 0.000380436190835 & 7.27740118397e-5 & 0.000658941961158 &
& {[}T3, row C 3{]} \\
S & 3 & 37 & 0.000775501992792 & 0.000343310022299 & 0.000939537300147 &
& {[}T3, row S 3{]} \\
X & 3 & 36 & 2.26586896653e-5 & 2.21983704619e-5 & 2.30966041916e-5 & &
{[}T3, row X 3{]} \\
S & 4 & 37 & 9.63721942695e-5 & 4.60435510795e-5 & 0.000139046731348 & &
{[}T3, row S 4{]} \\
X & 4 & 37 & 2.32406707119e-5 & 2.32155590409e-5 & 2.32870005206e-5 & &
{[}T3, row X 4{]} \\
S & 5 & 1 & 6.01344714582e-5 & 6.01344714582e-5 & 6.01344714582e-5 & not
decidable & {[}T3, row S 5{]} \\
X & 5 & 40 & 1.92365473742e-5 & 1.92334364212e-5 & 1.92413315356e-5 & &
{[}T3, row X 5{]} \\
X & 6 & 40 & 1.6241799957e-5 & 1.62397348555e-5 & 1.62425706568e-5 & &
{[}T3, row X 6{]} \\
\end{longtable}
}

No cell median of R(n\_H2) reaches the threshold of 0.10. The largest
cell median is 0.000775501992792 (class S, K index 3). The direction of
the NEGATIVE verdict therefore holds in every cell as well as in the
pooled set. This is exploratory item D7 (Section 9) and is not a
decision rule. The cell of class S at K index 5 has 1 orbit, below the
minimum set size of 20. It is not decidable and is neither pass nor
fail. The bootstrap resamples orbits as if they were independent draws
(limitation 1, Section 10). The cell intervals of Table 6 are
descriptive. Figure F2 shows R(n) by class and K, with the pooled series
over W.

\subsection{5. Results: H3 (claim C2),
SUPPORTED}\label{results-h3-claim-c2-supported}

\subsubsection{5.1 Preregistered rule}\label{preregistered-rule-2}

n\_conv is the smallest n after which every method stays within a 0.05
relative band of lambda\_ref up to N. Per orbit it is the maximum over
the methods of V. H3 is SUPPORTED if the ratio of the median n\_conv
over W to the median n\_conv over the strong-chaos set S is at least 10,
and the lower end of the 95 percent bootstrap interval of the ratio is
above 3. Here W is the set of Section 2.2 and S is the class C orbits at
K = 2 pi and K = 10. Orbits whose n\_conv is undefined are given n\_conv
= N, and their share is reported.

\subsubsection{5.2 Measured}\label{measured-2}

Table 7 gives the H3 measured values.

Table 7. H3 measured values.

{\def\LTcaptype{none} % do not increment counter
\begin{longtable}[]{@{}
  >{\raggedright\arraybackslash}p{(\linewidth - 6\tabcolsep) * \real{0.2500}}
  >{\raggedright\arraybackslash}p{(\linewidth - 6\tabcolsep) * \real{0.2500}}
  >{\raggedright\arraybackslash}p{(\linewidth - 6\tabcolsep) * \real{0.2500}}
  >{\raggedright\arraybackslash}p{(\linewidth - 6\tabcolsep) * \real{0.2500}}@{}}
\toprule\noalign{}
\begin{minipage}[b]{\linewidth}\raggedright
Quantity
\end{minipage} & \begin{minipage}[b]{\linewidth}\raggedright
Value
\end{minipage} & \begin{minipage}[b]{\linewidth}\raggedright
Threshold
\end{minipage} & \begin{minipage}[b]{\linewidth}\raggedright
Pointer
\end{minipage} \\
\midrule\noalign{}
\endhead
\bottomrule\noalign{}
\endlastfoot
Orbits in W & 355 & & {[}T4, header line{]} \\
Orbits in S & 80 & & {[}T4, header line{]} \\
Median n\_conv over W & 2.511886e6 & & {[}T4, header line{]} \\
Median n\_conv over S & 224.0 & & {[}T4, header line{]} \\
Ratio & 11213.7767857 & at least 10 & {[}T4, header line; params:
h3\_ratio{]} \\
95 percent bootstrap interval of the ratio & {[}7924.465, 14117.3125{]}
& lower end above 3 & {[}T4, header line; params: h3\_interval{]} \\
Share of undefined n\_conv & 0.0229885057471 & & {[}T4, header
line{]} \\
S orbits without positive growth & 0 & & {[}T4, header line{]} \\
\end{longtable}
}

H3 is supported. The verdict as printed is supported. Both conditions
pass. The ratio is at least 10 and the lower end is above 3. The
strong-chaos set S of H3 is the class C orbits at K = 2 pi and K = 10.
It is not the panel class S. The orbits with undefined n\_conv enter the
medians at the value N. Table 8 gives n\_conv by cell. The rows of class
C at K index 5 and 6 are the orbits of S. The other rows match the cells
of Table 6 in class, K index and orbit count. The cell-level medians are
descriptive. The class S cell at K index 5 has 1 orbit and is not
decidable, as in H2.

Table 8. H3: n\_conv by cell (descriptive).

{\def\LTcaptype{none} % do not increment counter
\begin{longtable}[]{@{}
  >{\raggedright\arraybackslash}p{(\linewidth - 14\tabcolsep) * \real{0.1250}}
  >{\raggedright\arraybackslash}p{(\linewidth - 14\tabcolsep) * \real{0.1250}}
  >{\raggedright\arraybackslash}p{(\linewidth - 14\tabcolsep) * \real{0.1250}}
  >{\raggedright\arraybackslash}p{(\linewidth - 14\tabcolsep) * \real{0.1250}}
  >{\raggedright\arraybackslash}p{(\linewidth - 14\tabcolsep) * \real{0.1250}}
  >{\raggedright\arraybackslash}p{(\linewidth - 14\tabcolsep) * \real{0.1250}}
  >{\raggedright\arraybackslash}p{(\linewidth - 14\tabcolsep) * \real{0.1250}}
  >{\raggedright\arraybackslash}p{(\linewidth - 14\tabcolsep) * \real{0.1250}}@{}}
\toprule\noalign{}
\begin{minipage}[b]{\linewidth}\raggedright
Class
\end{minipage} & \begin{minipage}[b]{\linewidth}\raggedright
K index
\end{minipage} & \begin{minipage}[b]{\linewidth}\raggedright
Orbits
\end{minipage} & \begin{minipage}[b]{\linewidth}\raggedright
Median
\end{minipage} & \begin{minipage}[b]{\linewidth}\raggedright
Lo
\end{minipage} & \begin{minipage}[b]{\linewidth}\raggedright
Hi
\end{minipage} & \begin{minipage}[b]{\linewidth}\raggedright
Note
\end{minipage} & \begin{minipage}[b]{\linewidth}\raggedright
Pointer
\end{minipage} \\
\midrule\noalign{}
\endhead
\bottomrule\noalign{}
\endlastfoot
C & 2 & 36 & 7.079458e6 & 6.309573e6 & 7.943282e6 & & {[}T4, row C
2{]} \\
S & 2 & 21 & 7.079458e6 & 5.0e6 & 7.079458e6 & & {[}T4, row S 2{]} \\
X & 2 & 36 & 7.079458e6 & 5.623413e6 & 7.079458e6 & & {[}T4, row X
2{]} \\
C & 3 & 34 & 4.223954e6 & 3.03001e6 & 6.309573e6 & & {[}T4, row C
3{]} \\
S & 3 & 37 & 5.623413e6 & 3.548134e6 & 7.943282e6 & & {[}T4, row S
3{]} \\
X & 3 & 36 & 5.623413e6 & 3.981072e6 & 6.6945155e6 & & {[}T4, row X
3{]} \\
S & 4 & 37 & 794328.0 & 707946.0 & 1.0e6 & & {[}T4, row S 4{]} \\
X & 4 & 37 & 63096.0 & 44668.0 & 89125.0 & & {[}T4, row X 4{]} \\
C & 5 & 40 & 266.5 & 200.0 & 335.5 & & {[}T4, row C 5{]} \\
S & 5 & 1 & 501187.0 & 501187.0 & 501187.0 & not decidable & {[}T4, row
S 5{]} \\
X & 5 & 40 & 251.0 & 168.0 & 447.0 & & {[}T4, row X 5{]} \\
C & 6 & 40 & 212.0 & 178.0 & 251.0 & & {[}T4, row C 6{]} \\
X & 6 & 40 & 158.0 & 133.5 & 224.0 & & {[}T4, row X 6{]} \\
\end{longtable}
}

Figure F3 shows the n\_conv distributions.

\subsection{\texorpdfstring{6. Results: H4 (REPRODUCES,
\citep{SalaArtuso2015}), DID NOT
REPRODUCE}{6. Results: H4 (REPRODUCES, {[}@SalaArtuso2015{]}), DID NOT REPRODUCE}}\label{results-h4-reproduces-salaartuso2015-did-not-reproduce}

\subsubsection{6.1 Preregistered rule}\label{preregistered-rule-3}

The start is x0 = 1e-3 and p0 = -1e-3, with T = 1e9 iterations. At K =
pi/2 (reference 0.298) both L1 and L5 (Float64, tau = 1) must be within
0.01. At K = 2 pi (reference 1.172) both must be within 0.02. The
reference values are those of Sala and Artuso \citep{SalaArtuso2015}. H4
is SUPPORTED if both K pass. A failure is reported as ``did not
reproduce'' with the numbers.

\subsubsection{6.2 Measured}\label{measured-3}

The verdict as printed is did not reproduce. Table 9 gives the verdict
table.

Table 9. H4 verdict table.

{\def\LTcaptype{none} % do not increment counter
\begin{longtable}[]{@{}
  >{\raggedright\arraybackslash}p{(\linewidth - 12\tabcolsep) * \real{0.1429}}
  >{\raggedright\arraybackslash}p{(\linewidth - 12\tabcolsep) * \real{0.1429}}
  >{\raggedright\arraybackslash}p{(\linewidth - 12\tabcolsep) * \real{0.1429}}
  >{\raggedright\arraybackslash}p{(\linewidth - 12\tabcolsep) * \real{0.1429}}
  >{\raggedright\arraybackslash}p{(\linewidth - 12\tabcolsep) * \real{0.1429}}
  >{\raggedright\arraybackslash}p{(\linewidth - 12\tabcolsep) * \real{0.1429}}
  >{\raggedright\arraybackslash}p{(\linewidth - 12\tabcolsep) * \real{0.1429}}@{}}
\toprule\noalign{}
\begin{minipage}[b]{\linewidth}\raggedright
Start
\end{minipage} & \begin{minipage}[b]{\linewidth}\raggedright
L1
\end{minipage} & \begin{minipage}[b]{\linewidth}\raggedright
L5
\end{minipage} & \begin{minipage}[b]{\linewidth}\raggedright
Reference
\end{minipage} & \begin{minipage}[b]{\linewidth}\raggedright
Tolerance
\end{minipage} & \begin{minipage}[b]{\linewidth}\raggedright
Pass
\end{minipage} & \begin{minipage}[b]{\linewidth}\raggedright
Pointer
\end{minipage} \\
\midrule\noalign{}
\endhead
\bottomrule\noalign{}
\endlastfoot
K = pi/2 & 0.256066843848 & 0.256066843848 & 0.298 & 0.01 & false &
{[}T5, row K = pi/2{]} \\
K = 2 pi & 1.17018723872 & 1.17018723872 & 1.172 & 0.02 & true & {[}T5,
row K = 2 pi{]} \\
\end{longtable}
}

H4 did not reproduce. At K = pi/2 both L1 and L5 gave 0.256066843848
against the reference 0.298. The value lies below the reference and
outside the tolerance of 0.01. The K = pi/2 start did not reproduce the
published value \citep{SalaArtuso2015}. At K = 2 pi the value
1.17018723872 against 1.172 is within the tolerance of 0.02. The
reproduction as a whole fails. The possibility that the orbit of the
published paper sits in a different component remains open. The tables
do not decide it. Each K rests on one start, with no error bar.

\subsubsection{6.3 Comparators (no
verdict)}\label{comparators-no-verdict}

Table 10 gives three comparators at the five K values: the ensemble mean
of lambda\_L5(N\_1) over the kept chaotic-sea orbits, the fit of
Shevchenko \citep{Shevchenko2014}, and the line ln(K/2)
\citep{Chirikov1979}. They never pass or fail. The frozen text
designates the line for K = pi/2, 2 pi and 10. The program prints it at
all five K. The rows at K = 0.5 and K = K\_c lie outside the designated
K and are shown as printed; the line takes negative values there
(-1.38629436112 and -0.721921822156). The range of applicability of the
fit and of the line is not checked in this paper. No verdict is attached
to any comparator row.

Table 10. Comparators.

{\def\LTcaptype{none} % do not increment counter
\begin{longtable}[]{@{}
  >{\raggedright\arraybackslash}p{(\linewidth - 12\tabcolsep) * \real{0.1429}}
  >{\raggedright\arraybackslash}p{(\linewidth - 12\tabcolsep) * \real{0.1429}}
  >{\raggedright\arraybackslash}p{(\linewidth - 12\tabcolsep) * \real{0.1429}}
  >{\raggedright\arraybackslash}p{(\linewidth - 12\tabcolsep) * \real{0.1429}}
  >{\raggedright\arraybackslash}p{(\linewidth - 12\tabcolsep) * \real{0.1429}}
  >{\raggedright\arraybackslash}p{(\linewidth - 12\tabcolsep) * \real{0.1429}}
  >{\raggedright\arraybackslash}p{(\linewidth - 12\tabcolsep) * \real{0.1429}}@{}}
\toprule\noalign{}
\begin{minipage}[b]{\linewidth}\raggedright
K index
\end{minipage} & \begin{minipage}[b]{\linewidth}\raggedright
K
\end{minipage} & \begin{minipage}[b]{\linewidth}\raggedright
Orbits
\end{minipage} & \begin{minipage}[b]{\linewidth}\raggedright
Mean
\end{minipage} & \begin{minipage}[b]{\linewidth}\raggedright
Shevchenko fit
\end{minipage} & \begin{minipage}[b]{\linewidth}\raggedright
ln(K/2)
\end{minipage} & \begin{minipage}[b]{\linewidth}\raggedright
Pointer
\end{minipage} \\
\midrule\noalign{}
\endhead
\bottomrule\noalign{}
\endlastfoot
2 & 0.5 & 40 & 0.0575923614667 & 0.06665 & -1.38629436112 & {[}T5,
comparators, K index 2{]} \\
3 & 0.971635406 & 40 & 0.091095497804 & 0.12951899962 & -0.721921822156
& {[}T5, comparators, K index 3{]} \\
4 & 1.57079632679 & 40 & 0.298004458001 & 0.342657808664 &
-0.24156447527 & {[}T5, comparators, K index 4{]} \\
5 & 6.28318530718 & 40 & 1.17214446011 & 1.17006018176 & 1.14472988585 &
{[}T5, comparators, K index 5{]} \\
6 & 10 & 40 & 1.62000488786 & 1.61943791243 & 1.60943791243 & {[}T5,
comparators, K index 6{]} \\
\end{longtable}
}

\subsubsection{6.4 Exploratory item D4
(EXPLORATORY)}\label{exploratory-item-d4-exploratory}

D4 (EXPLORATORY): H4 at the shorter horizon. The cross-check uses T =
1e8. It has no verdict under the frozen rule. D4 does not change the H4
verdict, which stands under the frozen rule: H4 did not reproduce. At K
= pi/2 it gives 0.299256534631 (Float64) and 0.299012763029 (Double64).
At K = 2 pi it gives 1.17243899416 and 1.17179404609. The Float64 value
at K = pi/2 at T = 1e9 in the verdict table is 0.256066843848. The
ensemble mean of lambda\_L5(N\_1) over the 40 kept chaotic-sea orbits at
K = pi/2 is 0.298004458001. The tables do not show why the single start
at K = pi/2 gives a different value at T = 1e9 than at T = 1e8. No cause
is claimed.

\subsection{7. Results: H5, SUPPORTED}\label{results-h5-supported}

\subsubsection{7.1 Preregistered rule}\label{preregistered-rule-4}

On the ladder subset S\_L, H5a is the fraction of S\_L orbits whose two
labels are identical at Float64, Double64 and BigFloat. It must be at
least 0.90. H5b is the fraction of all panel orbits with either axis
``unresolved''. It must be at most 0.30. H5c is the fraction with
structure ``unresolved'' among orbits with instability ``none detected
on the tested horizon''. It must be at most 0.30. H5 is supported only
if all three sub-rules pass. H5a is read from the per-precision labels
before the ladder test acts.

\subsubsection{7.2 Measured}\label{measured-4}

The verdict as printed is supported. Table 11 gives the three sub-rules.

Table 11. H5 sub-rules.

{\def\LTcaptype{none} % do not increment counter
\begin{longtable}[]{@{}
  >{\raggedright\arraybackslash}p{(\linewidth - 10\tabcolsep) * \real{0.1667}}
  >{\raggedright\arraybackslash}p{(\linewidth - 10\tabcolsep) * \real{0.1667}}
  >{\raggedright\arraybackslash}p{(\linewidth - 10\tabcolsep) * \real{0.1667}}
  >{\raggedright\arraybackslash}p{(\linewidth - 10\tabcolsep) * \real{0.1667}}
  >{\raggedright\arraybackslash}p{(\linewidth - 10\tabcolsep) * \real{0.1667}}
  >{\raggedright\arraybackslash}p{(\linewidth - 10\tabcolsep) * \real{0.1667}}@{}}
\toprule\noalign{}
\begin{minipage}[b]{\linewidth}\raggedright
Sub-rule
\end{minipage} & \begin{minipage}[b]{\linewidth}\raggedright
Orbits
\end{minipage} & \begin{minipage}[b]{\linewidth}\raggedright
Fraction
\end{minipage} & \begin{minipage}[b]{\linewidth}\raggedright
Limit
\end{minipage} & \begin{minipage}[b]{\linewidth}\raggedright
Result
\end{minipage} & \begin{minipage}[b]{\linewidth}\raggedright
Pointer
\end{minipage} \\
\midrule\noalign{}
\endhead
\bottomrule\noalign{}
\endlastfoot
H5a & 92 & 1.0 & at least 0.90 & pass & {[}T7, row H5a; params:
h5a\_min{]} \\
H5b & 883 & 0.0951302378256 & at most 0.30 & pass & {[}T7, row H5b;
params: h5b\_max{]} \\
H5c & 334 & 0.0568862275449 & at most 0.30 & pass & {[}T7, row H5c;
params: h5c\_max{]} \\
\end{longtable}
}

All three label conditions are met. H5 is supported under its three
sub-rules.

\subsubsection{7.3 Exploratory items D3 and D5
(EXPLORATORY)}\label{exploratory-items-d3-and-d5-exploratory}

D3 (EXPLORATORY): H5a is 1.0 while many ladder labels are forced. H5a is
1.0 on 92 S\_L orbits, read before the ladder test acts. The ladder test
then forced the labels of 35 of those 92 orbits to unresolved, a share
of 0.380434782609. Both hold at once: the per-label agreement across
precisions, and the override of 35 labels by the ladder test. The
program prints the unresolved share before the tests as 46 of 883
(0.0520951302378) and after as 84 of 883 (0.0951302378256). The second
value is the H5b fraction. Other forced tests: direction 7 of 172
(0.0406976744186). Estimator, diagnostics and excluded show 0 of 883
each. The tables do not show how the forced counts (35 and 7) relate to
the rise from 46 to 84. No cause is claimed.

D5 (EXPLORATORY): unresolved labels by cell. ``Both'' counts orbits with
both axes unresolved (65 in all). ``Structure only'' counts orbits with
instability ``none detected on the tested horizon'' and structure
unresolved (19 in all). The two add to the 84 of H5b. Cells not listed
in Table 12 have none.

Table 12. Unresolved labels by cell (EXPLORATORY).

{\def\LTcaptype{none} % do not increment counter
\begin{longtable}[]{@{}
  >{\raggedright\arraybackslash}p{(\linewidth - 10\tabcolsep) * \real{0.1667}}
  >{\raggedright\arraybackslash}p{(\linewidth - 10\tabcolsep) * \real{0.1667}}
  >{\raggedright\arraybackslash}p{(\linewidth - 10\tabcolsep) * \real{0.1667}}
  >{\raggedright\arraybackslash}p{(\linewidth - 10\tabcolsep) * \real{0.1667}}
  >{\raggedright\arraybackslash}p{(\linewidth - 10\tabcolsep) * \real{0.1667}}
  >{\raggedright\arraybackslash}p{(\linewidth - 10\tabcolsep) * \real{0.1667}}@{}}
\toprule\noalign{}
\begin{minipage}[b]{\linewidth}\raggedright
K index
\end{minipage} & \begin{minipage}[b]{\linewidth}\raggedright
Class
\end{minipage} & \begin{minipage}[b]{\linewidth}\raggedright
Cell orbits
\end{minipage} & \begin{minipage}[b]{\linewidth}\raggedright
Both
\end{minipage} & \begin{minipage}[b]{\linewidth}\raggedright
Structure only
\end{minipage} & \begin{minipage}[b]{\linewidth}\raggedright
Pointer
\end{minipage} \\
\midrule\noalign{}
\endhead
\bottomrule\noalign{}
\endlastfoot
0 & controls & 13 & 3 & 0 & {[}T1, row 0 controls; T6, rows 0
controls{]} \\
2 & C & 40 & 4 & 0 & {[}T1, row 2 C; T6, rows 2 C{]} \\
2 & R & 40 & 5 & 0 & {[}T1, row 2 R; T6, rows 2 R{]} \\
2 & S & 25 & 4 & 0 & {[}T1, row 2 S; T6, rows 2 S{]} \\
2 & T & 40 & 0 & 11 & {[}T1, row 2 T; T6, rows 2 T{]} \\
2 & X & 40 & 4 & 0 & {[}T1, row 2 X; T6, rows 2 X{]} \\
3 & C & 40 & 6 & 0 & {[}T1, row 3 C; T6, rows 3 C{]} \\
3 & R & 40 & 6 & 0 & {[}T1, row 3 R; T6, rows 3 R{]} \\
3 & S & 40 & 3 & 0 & {[}T1, row 3 S; T6, rows 3 S{]} \\
3 & T & 40 & 2 & 5 & {[}T1, row 3 T; T6, rows 3 T{]} \\
3 & X & 40 & 4 & 0 & {[}T1, row 3 X; T6, rows 3 X{]} \\
4 & C & 40 & 2 & 0 & {[}T1, row 4 C; T6, rows 4 C{]} \\
4 & R & 40 & 6 & 1 & {[}T1, row 4 R; T6, rows 4 R{]} \\
4 & S & 40 & 3 & 0 & {[}T1, row 4 S; T6, rows 4 S{]} \\
4 & T & 40 & 2 & 2 & {[}T1, row 4 T; T6, rows 4 T{]} \\
4 & X & 40 & 3 & 0 & {[}T1, row 4 X; T6, rows 4 X{]} \\
5 & S & 5 & 4 & 0 & {[}T1, row 5 S; T6, rows 5 S{]} \\
5 & T & 40 & 4 & 0 & {[}T1, row 5 T; T6, rows 5 T{]} \\
\end{longtable}
}

K index 1, K index 6 and K index 5 classes C and X have no unresolved
orbit.

\subsection{8. Cost and sensitivity (reported without a decision
rule)}\label{cost-and-sensitivity-reported-without-a-decision-rule}

\subsubsection{8.1 Cost}\label{cost}

Table 13 gives ns per orbit-step by method and precision, and relative
to L5 in the same precision. Each row is an unweighted mean over all
items of that method and precision, so the item sets differ between
rows. Relative to L5 in Float64 (168.60570837 ns per orbit-step), L1
costs 0.701549448499, L2 0.503966109187, L4 0.982175001597, L9
1.23876330018 and L6 2.10273197634. These ratios compare rows with
different item sets (1917, 1571, 1915, 1915 and 1227 items against 2605
for L5), and they come from a single run that used 23 of 23 available
threads, with no repeated timing reported. L5 costs 4882.6924383 ns in
Double64 and 377399.856865 ns in BigFloat(256). The Double64 and
BigFloat rows have 186 and 92 items against 2605 for L5 in Float64, so
they are also not like-for-like with the Float64 rows.

Table 13. Cost per orbit-step.

{\def\LTcaptype{none} % do not increment counter
\begin{longtable}[]{@{}
  >{\raggedright\arraybackslash}p{(\linewidth - 10\tabcolsep) * \real{0.1667}}
  >{\raggedright\arraybackslash}p{(\linewidth - 10\tabcolsep) * \real{0.1667}}
  >{\raggedright\arraybackslash}p{(\linewidth - 10\tabcolsep) * \real{0.1667}}
  >{\raggedright\arraybackslash}p{(\linewidth - 10\tabcolsep) * \real{0.1667}}
  >{\raggedright\arraybackslash}p{(\linewidth - 10\tabcolsep) * \real{0.1667}}
  >{\raggedright\arraybackslash}p{(\linewidth - 10\tabcolsep) * \real{0.1667}}@{}}
\toprule\noalign{}
\begin{minipage}[b]{\linewidth}\raggedright
Method
\end{minipage} & \begin{minipage}[b]{\linewidth}\raggedright
Precision
\end{minipage} & \begin{minipage}[b]{\linewidth}\raggedright
Items
\end{minipage} & \begin{minipage}[b]{\linewidth}\raggedright
ns per orbit-step
\end{minipage} & \begin{minipage}[b]{\linewidth}\raggedright
Relative to L5
\end{minipage} & \begin{minipage}[b]{\linewidth}\raggedright
Pointer
\end{minipage} \\
\midrule\noalign{}
\endhead
\bottomrule\noalign{}
\endlastfoot
L5 & Float64 & 2605 & 168.60570837 & 1.0 & {[}cost section, row L5
Float64{]} \\
L1 & Float64 & 1917 & 118.285241721 & 0.701549448499 & {[}cost section,
row L1 Float64{]} \\
L2 & Float64 & 1571 & 84.9715628342 & 0.503966109187 & {[}cost section,
row L2 Float64{]} \\
L4 & Float64 & 1915 & 165.600311888 & 0.982175001597 & {[}cost section,
row L4 Float64{]} \\
L9 & Float64 & 1915 & 208.86256373 & 1.23876330018 & {[}cost section,
row L9 Float64{]} \\
L6 & Float64 & 1227 & 354.532614384 & 2.10273197634 & {[}cost section,
row L6 Float64{]} \\
L5 & Double64 & 186 & 4882.6924383 & 1.0 & {[}cost section, row L5
Double64{]} \\
L5 & BigFloat(256) & 92 & 377399.856865 & 1.0 & {[}cost section, row L5
BigFloat(256){]} \\
\end{longtable}
}

The item count of the L1 Float64 row, 1917, is from the table row. The
findings file lists that row without a count. Run totals: the total wall
time was 11957.993616580963 seconds against a cap of 12.0 hours. The
plan had 1363 chunks and 15569 items. The run used 23 threads of 23
available. All 13 arms are complete.

\subsubsection{8.2 Sensitivity}\label{sensitivity}

The sensitivity table has 31 rows. The column is the median paired
difference of lambda at n\_H2 against the primary item, with 171 orbits
per row except the tau rows. These rows are evaluated at n\_H2 = 1e4.
H1c is evaluated at N\_1 = 1e7. The rows below are therefore not an
evaluation of the H1c result.

Direction. L1 and L5 values are of order 1e-9 at offset 0.3 (L5
1.73141166638e-9). They are of order 1e-5 or 1e-6 at the other offsets:
L5 is -3.32612943661e-5 at 0.3 + 2pi/5, -3.23438506002e-5 at 0.3 + 3pi/5
and -3.48457834277e-6 at 0.3 + pi/5. L9 is 0.0 in every direction row.

The eps rows for L2. With eps = 1e-10 the value is 2.26005874629e-9.
With eps = 1e-12 it is 5.77729472453e-5. With eps = 1e-6 it is
4.04116577712e-10.

Wrapped versus lifted is 0.0 for all eight methods of V. The column
median\_log10\_nconv\_ratio is 0.0 in all 31 rows. This is reported as
printed. It has not been checked whether a constant column is an
artefact.

Tau sweep. Tau = 2 and 4 show no breakdown, with 171 orbits for L4 and
L5. Tau = 8 has 154 orbits for L4 and L5, with breakdown count 17. Tau =
16 has 132 orbits for L4 with breakdown count 39 and 134 orbits for L5
with breakdown count 37. At tau = 16 the breakdown steps are 16 or 32
for L4 and 16 to 288 for L5.

Symplectic residuals eps\_pair(n) and eps\_det(n) are shown in Figure F5
only. No numbers for them are quoted. Figure F1 shows the lambda\_m(n)
curves.

\subsection{9. Exploratory results}\label{exploratory-results}

All items in this section are EXPLORATORY. They are not confirmatory,
have no frozen rule and carry no claimed cause. Items D1 to D5 appear in
Sections 3, 6 and 7, where they bear on a hypothesis. D6 and D7 are
reported here.

D6 (EXPLORATORY): short, empty and not decidable cells. The panel
manifest marks two cells as short (``fewer qualifying grid cells than
n\_c''): K index 2 class S (25 drawn, 25 qualifying) and K index 5 class
S (5 drawn, 5 qualifying). It marks as empty K index 5 class R (no
elliptic periodic point) and K index 6 classes R (no elliptic periodic
point), T and S (no qualifying grid cell). It marks K index 1 classes C,
S and X as undefined (class undefined at K = 0). Class X shows 0
qualifying grid cells at every K index from 2 to 6. The preregistered
selection rule does not draw class X from the screening and places the
starts near the hyperbolic fixed point. In the H2 and H3 tables only
class S at K index 5 is not decidable, with 1 orbit.

D7 (EXPLORATORY): cell patterns. The largest cell median of R(n\_H2) is
0.000775501992792 (class S, K index 3). The smallest is
5.82402750446e-12 (class X, K index 2). Both are in Table 6. The cell
medians of n\_conv are between 4.223954e6 and 7.079458e6 at K index 2
and 3. The horizon there is N\_1 = 1e7. These values are in Table 8. At
K index 4 they are 794328.0 (class S) and 63096.0 (class X). At K index
5 and 6 they lie between 158.0 and 266.5, excluding class S at K index
5. Several cells share the same median (for example 7.079458e6 in three
cells of Table 8), which is consistent with n\_conv taking recorded
values of n (Section 2.2). The effect of this resolution on the
bootstrap intervals is not examined. The tables do not show why.

\subsection{10. Limitations}\label{limitations}

\begin{enumerate}
\def\labelenumi{\arabic{enumi}.}
\tightlist
\item
  The study rests on one draw of the panel (a single seed) with 883
  orbits. The bootstrap intervals resample orbits within that draw,
  treat the orbits as independent resampling units, and do not reflect a
  second draw. The cell tables of H2 and H3 are descriptive. The
  composition of the 883 orbits is given in Section 2.3.
\item
  The panel is run in Float64 only. Precision comparisons come from the
  ladder subset S\_L of 92 orbits, and H5a rests on that subset.
\item
  A set or cell with fewer than 20 orbits is not decidable and is
  neither pass nor fail. In the H2 and H3 tables only class S at K index
  5 (5 drawn, 1 in W) is below that minimum. Class S at K index 2 has 25
  drawn and 21 in W, which is above the minimum. Cell-level medians and
  intervals are descriptive, and the decisions use the pooled sets. H4
  rests on one start per K, as the frozen text specifies, with no error
  bar.
\item
  L2 runs in wrapped coordinates. This affects the H1c metric, the
  sub-rule that failed. The tables do not show whether this accounts for
  any part of the H1c result.
\item
  The numbers were checked against the tables file of the run. The
  analysis program was verified against an independent recomputation
  (\texttt{analysis-verification.md}). The numbers were not recomputed
  from the stored series for the results document. Figures F1 to F7 were
  redrawn and viewed at paper stage 04; they plot the figure data of the
  tables file, unchanged.
\item
  W is defined here by the preregistration, as in the results document
  and the findings file. The analysis program still prints a shorter
  phrase for W in its list of readings (class C at 0 \textless{} K
  \textless= K\_c), which does not describe the cells of Tables 6 and 8;
  the membership of W (355 orbits) does not depend on that phrase.
\item
  The controls table lists 10 of the 13 controls; the other 3 are the
  elliptic controls with exponent 0, which no decision rule of H1 uses
  (Section 3.3). Orbit 881 is an info control and the orbit dropped as
  drift.
\item
  Table 1 gives two git commits, one for the run and one for the
  analysis. They are related through the repository history only; the
  study manifests record each commit separately and no single run record
  names both.
\item
  The repository folder \texttt{data/} holds run outputs and is
  git-ignored (project folder map), so it is not part of the repository
  at the URL of Section 11. The stored series are packed for deposit as
  a dataset on Zenodo (1363 chunk files with the manifest, the panel,
  the preregistration, the parameter file, the file-format description
  and a SHA-256 list); its DOI is 10.5281/zenodo.23534755 (Section 11).
\item
  The frozen preregistration is in the repository (Section 11). The
  results document records no external time-stamped deposit of it.
\item
  Cost values come from a single run in which 23 of 23 available threads
  were used, with no repeated timing reported. The item sets differ
  between rows, and the Double64 and BigFloat rows have 186 and 92
  items. The cost ratios are therefore not controlled comparisons.
\item
  n\_conv takes recorded values of n, because the series are stored at
  20 points per decade, log-spaced, plus the first 100 steps. Its
  resolution and its effect on the bootstrap intervals are not examined.
\item
  The column median\_log10\_nconv\_ratio is 0.0 in all 31 sensitivity
  rows. It has not been checked whether this is an artefact.
\item
  Open items carried over from the study hand-off: the STRONG\_K
  constant was kept by owner decision, and its role is not described in
  the results document. The package requires Julia 1.11 or later, so no
  dependency manifest for Julia 1.10 is kept.
\end{enumerate}

\subsection{11. Data and code
availability}\label{data-and-code-availability}

The repository is https://github.com/georgehadji/Chaos
(GPL-3.0-or-later). The frozen preregistration is
\texttt{workspaces/study/stages/01-design/output/p1-le-benchmark-preregistration.md}
in the repository. Its header SHA-256 is in Table 1. The git SHA at
freeze recorded in its header is
\seqsplit{1e760e229eabf9d8323ad4d8097cc1dbd017f845}, the parent commit of the
freeze commit. Amendments A1 to A6 are in the same file and in its
amendment log. The tagged release v0.1.0 was cut at milestone M3 exit,
at commit 238e96b, a descendant of both commits of Table 1. The run used
package version 0.0.1 at those commits, before the tag. The tagged
release v0.1.1, at commit efeac86, adds the user documentation; between
the analysis commit and v0.1.1 the package code changes only in
docstrings and in a file-descriptor conversion of the chunk writer that
is a no-op on Julia 1.13.2, the Julia version of the run. The run
manifest and the stored series identify the run. The SHA-256 values are
in Table 1. The stored series are deposited as a Zenodo dataset
(limitation 9, Section 10); dataset DOI:
\href{https://doi.org/10.5281/zenodo.23534755}{10.5281/zenodo.23534755}.

\subsubsection{Availability statement}\label{availability-statement}

Repository: \url{https://github.com/georgehadji/Chaos}
(GPL-3.0-or-later). The numbers in this paper come from one analysis run
at commit \seqsplit{d31b655d0c33e2f2f1a67bbbdeba9c8c8a3fd036}. The stored series
come from the run at commit \seqsplit{3e405f3324f8d25162580d3380db26c877d7ca63}, an
ancestor of the analysis commit with the same series code (Section 2.4).
The tagged release v0.1.0 is at commit 238e96b; the run used package
version 0.0.1, before the tag. The release to cite for the code is
v0.1.1 (commit efeac86), archived with \texttt{CITATION.cff}. Study
manifest:
\texttt{workspaces/study/stages/03-run/output/p1-le-benchmark-manifest.json}
(SHA-256
\seqsplit{ae6ec80eb08fbd6e98aeef1fbf840a567a03102ab943f2ae0aaf2b3e367f191d}).
Results document:
\texttt{workspaces/study/stages/05-report/output/p1-le-benchmark-results.md}.
The stored series are not part of the repository; they are deposited as
a Zenodo dataset (DOI 10.5281/zenodo.23534755).

\subsection{AI-use disclosure}\label{ai-use-disclosure}

Drafting, adversarial review and revision of this manuscript, and the
analysis code, were done with Claude models (Anthropic, Claude Sonnet
5.5) operated through Claude Code under the direction of the author. The
author decided the scope, the framing and the preregistered rules, and
accepted or redirected the output at each stage checkpoint. Numbers in
the text were copied from a results document produced by the analysis
program, which was checked against an independent recomputation.
Citations are limited to entries verified in the project bibliography.
The author is responsible for the content.

\subsection{12. Conclusion}\label{conclusion}

Scope: one panel from a single draw (a single seed), run in Float64,
with five nonzero kick values and one start per kick value in H4
(Section 10). Under the preregistered rules, H1 (claim C1) is falsified,
on sub-rule H1c alone, with H1a and H1b met in full. It was evaluated in
Float64 at N\_1 = 1e7, with L2 in wrapped coordinates (limitation 4).
Claim C2 has two preregistered tests with different verdicts. H2 (claim
C2) is not supported (NEGATIVE). It measures R(n\_H2) = Delta\_V(n\_H2)
/ lambda\_ref over the set W at n\_H2 = 1e4 steps. H3 (claim C2) is
supported. It compares the median n\_conv over W with the median n\_conv
over the strong-chaos set S, where n\_conv is the maximum over V of the
smallest n after which a method stays within a 0.05 relative band of
lambda\_ref. This paper does not combine the two verdicts into one
verdict on C2. H4 did not reproduce the published value at K = pi/2 and
passed at K = 2 pi. H5 is supported under its three sub-rules (H5a, H5b,
H5c). The ladder test forced 35 of the 92 ladder-subset labels to
unresolved (exploratory item D3). H4 and H5 are not tied to C1 or C2 in
the findings file. TOOL: Chaoslab v0.1.1 is a tagged and documented
release of the code the run used; every method of the run is
\texttt{verified} (Section 2.1, Section 11).

\subsection{Figure list}\label{figure-list}

The figure files are in the \texttt{figures/} folder next to this
manuscript. They were drawn from the figure data of the study's tables
file by the figure extension of the repository at commit caaf95a; the
figure data are those of the study output, unchanged. In this manuscript
an F-label always names a figure. The figures are listed here and not
embedded. The captions below describe what each figure draws.

\begin{itemize}
\tightlist
\item
  Figure F1: \texttt{p1-le-benchmark-fig-F1.pdf}. lambda\_1(n) of the
  nine methods for the first drawn orbit of each populated cell (design
  class and K index), with the band between the smallest and the largest
  of the eight methods of V (L2 excluded).
\item
  Figure F2: \texttt{p1-le-benchmark-fig-F2.pdf}. Median of R(n) over
  the orbits of each cell and over the pooled set W, with the 95 percent
  bootstrap band.
\item
  Figure F3: \texttt{p1-le-benchmark-fig-F3.pdf}. Empirical distribution
  functions of n\_conv per cell; colour by design class, line style by K
  index.
\item
  Figure F4: \texttt{p1-le-benchmark-fig-F4.pdf}. Share of orbits forced
  to unresolved by each label test, among the orbits the test applies
  to, and the share of all orbits unresolved before and after the tests.
\item
  Figure F5: \texttt{p1-le-benchmark-fig-F5.pdf}. Largest running
  maximum over the orbits of each stored diagnostic (det, symp, orth,
  pair, sum) per method, with the median threshold over the orbits; each
  orbit is tested against its own threshold.
\item
  Figure F6: \texttt{p1-le-benchmark-fig-F6.pdf}. Thread nanoseconds per
  orbit-step by method and precision.
\item
  Figure F7: \texttt{p1-le-benchmark-fig-F7.pdf}. Median paired
  difference of lambda(n\_H2) from the primary item for the levels of
  tau, eps, initial direction and coordinate mode, over the kept S\_D
  orbits.
\end{itemize}

\bibliography{refs.bib}

\end{document}
